% META: {"id":"1SPE-DERGLOBAL-CO-049","chapitre":"1SPE-DERIVATION-GLOBAL","type_objet":"corrige","exercice_id":"1SPE-DERGLOBAL-EX-049","status":"generated"}
% BEGIN-VERIFY
% from sympy import symbols, diff, solve, Rational
% x = symbols('x')
% P = -Rational(5,4)*x**2 + 45*x - 20
% Pp = diff(P, x)
% assert Pp == -Rational(5,2)*x + 45
% xopt = solve(Pp, x)
% assert xopt == [18]
% assert P.subs(x, 18) == 385
% Ppp = diff(Pp, x)
% assert Ppp == -Rational(5,2)
% assert Ppp < 0
% # Check sign of Pp
% assert Pp.subs(x, 10) > 0
% assert Pp.subs(x, 20) < 0
% END-VERIFY
\begin{corrige}{1SPE-DERGLOBAL-EX-049}
\textbf{1.} La recette est $R(x) = x \cdot p(x) = x(50 - x) = 50x - x^2$.

\textbf{2.} Le profit vaut :
\begin{align*}
P(x) &= R(x) - C(x) = 50x - x^2 - \frac{x^2}{4} - 5x - 20 \\
     &= -x^2 - \frac{x^2}{4} + 45x - 20 = -\frac{5x^2}{4} + 45x - 20.
\end{align*}

\textbf{3.} On dérive :
\[
P'(x) = -\frac{5}{2}x + 45.
\]
$P'(x) = 0 \iff x = \dfrac{45 \times 2}{5} = 18$.

$P'(x) > 0$ pour $x < 18$ et $P'(x) < 0$ pour $x > 18$.

\textbf{4.} Tableau de variations de $P$ sur $\left[0\,;\,40\right]$ :
\[
\begin{array}{|c|ccccc|}
\hline
x     & 0  &  & 18  &  & 40 \\
\hline
P'(x) &    & + & 0  & - & \\
\hline
P(x)  & -20 & \nearrow & 385 & \searrow & -300 \\
\hline
\end{array}
\]
($P(0) = -20$, $P(40) = -\frac{5 \times 1600}{4} + 45 \times 40 - 20 = -2000 + 1800 - 20 = -220$.)

\textbf{5.} Le profit est maximal pour $x^* = \boxed{18}$ unités.

Le profit maximal est :
\[
P(18) = -\frac{5 \times 324}{4} + 45 \times 18 - 20 = -405 + 810 - 20 = \boxed{385 \text{ euros}}.
\]

On vérifie : $P''(x) = -\dfrac{5}{2} < 0$ pour tout $x$, donc en particulier en $x^* = 18$,
ce qui confirme que c'est bien un maximum (la parabole est tournée vers le bas).
\end{corrige}
